% GENERATED FILE. Edit canonical lessons, then run npm run generate:books.
% canonical-source-hash: fnv1a64:86a16b05fd5a4cf2
% canonical-lessons: HI-C305-uphanna, HI-C305-badh, HI-C305-kramsah, HI-C305-sailab, HI-C305-bhavar

\chapter{The Water Rose}
\label{ch:hi-the-water-rose}
\hlchaptermodality{305}

\begin{chapteropening}
\textbf{By the end of this chapter:} \emph{I can say how a flood rose little by little.}

Complete the last lesson of chapter 305: I can say how a flood rose little by little.
\end{chapteropening}

\section[\texorpdfstring{\hi{उफनना} (uphannā)}{उफनना (uphannā)}]{\hi{उफनना} (uphannā) — to swell up; to boil over}

\label{lesson:HI-C305-uphanna}



\begin{quote}
Before the new one: say the Hindi for torrential (of rain), then the Hindi for destruction.
\end{quote}

\subsection*{You'll want to know: \hi{उफनना}}
\textbf{\hi{उफनना}} — \emph{uphannā} — \textquotedblleft{}to swell up; to boil over\textquotedblright{}.

\begin{quote}
\hi{भारी} \hi{बारिश} \hi{के} \hi{बाद} \hi{नदी} \hi{उफन} \hi{गई।}
\end{quote}

\begin{quote}
\emph{bhārī bāriś ke bād nadī uphan gaī.}
\end{quote}

\begin{quote}
After the heavy rain the river swelled over its banks.
\end{quote}

Past (it …): \hi{वह} \hi{नदी} \hi{उफनी}, \emph{vah nadī uphnī}.

\begin{cousinweb}
\textbf{\hi{उफनना}} (\emph{uphannā}) is inherited from \emph{*utphaṇati}, a form CDIAL reconstructs (CDIAL 1836). Milk on the fire \textbf{\hi{उफनता} \hi{है}}, and so does a river in flood.
\end{cousinweb}

\begin{grammarlens}[title={What you've built}]
One more word for telling what happened.
\end{grammarlens}

\subsection*{Your turn}
\begin{itemize}
\raggedright
  \item \emph{Say it:} \emph{uphannā}
  \item \emph{Say it:} \emph{uphannā}, once more
  \item \emph{Say it:} \emph{bhārī bāriś ke bād nadī uphan gaī.} — After the heavy rain the river swelled over its banks.
  \item \emph{Recall:} say the Hindi for torrential (of rain), then the Hindi for destruction, then say \emph{uphannā} again
\end{itemize}

\begin{tcolorbox}[breakable,colback=teal!4,colframe=teal!35!black,title={Before you move on}]
What does \hi{उफनना} mean? (\textquotedblleft{}To swell up; to boil over\textquotedblright{}.) Say it once more.
\end{tcolorbox}
\section[\texorpdfstring{\hi{बाँध} (bā̃dh)}{बाँध (bā̃dh)}]{\hi{बाँध} (bā̃dh) — a dam; an embankment}

\label{lesson:HI-C305-badh}



\begin{quote}
Before the new one: say the Hindi for destruction, then the Hindi for to swell up.
\end{quote}

\subsection*{You'll want to know: \hi{बाँध}}
\textbf{\hi{बाँध}} — \emph{bā̃dh} — \textquotedblleft{}a dam; an embankment\textquotedblright{}.

\begin{quote}
\hi{पानी} \hi{के} \hi{ज़ोर} \hi{से} \hi{बाँध} \hi{टूट} \hi{गया।}
\end{quote}

\begin{quote}
\emph{pānī ke zor se bā̃dh ṭūṭ gayā.}
\end{quote}

\begin{quote}
The force of the water broke the dam.
\end{quote}

Masculine.

\begin{cousinweb}
comes from Sanskrit \emph{bandha}, \textquotedblleft{}a binding, a fastening\textquotedblright{} (Monier-Williams), the family of \textbf{\hi{बाँधना}}, to tie: a \textbf{\hi{बाँध}} ties the water in.
\end{cousinweb}

\begin{grammarlens}[title={What you've built}]
One more word for telling what happened.
\end{grammarlens}

\subsection*{Your turn}
\begin{itemize}
\raggedright
  \item \emph{Say it:} \emph{bā̃dh}
  \item \emph{Say it:} \emph{bā̃dh}, once more
  \item \emph{Say it:} \emph{pānī ke zor se bā̃dh ṭūṭ gayā.} — The force of the water broke the dam.
  \item \emph{Recall:} say the Hindi for destruction, then the Hindi for to swell up, then say \emph{bā̃dh} again
\end{itemize}

\begin{tcolorbox}[breakable,colback=teal!4,colframe=teal!35!black,title={Before you move on}]
What does \hi{बाँध} mean? (\textquotedblleft{}A dam; an embankment\textquotedblright{}.) Say it once more.
\end{tcolorbox}
\section[\texorpdfstring{\hi{क्रमशः} (kramśaḥ)}{क्रमशः (kramśaḥ)}]{\hi{क्रमशः} (kramśaḥ) — little by little; by degrees}

\label{lesson:HI-C305-kramsah}



\begin{quote}
Before the new one: say the Hindi for to swell up, then the Hindi for a dam.
\end{quote}

\subsection*{You'll want to know: \hi{क्रमशः}}
\textbf{\hi{क्रमशः}} — \emph{kramśaḥ} — \textquotedblleft{}little by little; by degrees\textquotedblright{}.

\begin{quote}
\hi{नदी} \hi{का} \hi{पानी} \hi{क्रमशः} \hi{ऊपर} \hi{आ} \hi{रहा} \hi{था।}
\end{quote}

\begin{quote}
\emph{nadī kā pānī kramśaḥ ūpar ā rahā thā.}
\end{quote}

\begin{quote}
The river water was rising little by little.
\end{quote}

Formal and written; speech says \hi{धीरे}-\hi{धीरे}.

\begin{cousinweb}
\textbf{\hi{क्रमशः}} (\emph{kramśaḥ}) is Sanskrit \emph{kramaśas}, \textquotedblleft{}step by step, by degrees\textquotedblright{} (Monier-Williams), from \emph{krama}, a step. It ends in the sign \textbf{\hi{ः}}. At the end of a newspaper story printed in parts it means \textquotedblleft{}to be continued\textquotedblright{}.
\end{cousinweb}

\begin{grammarlens}[title={What you've built}]
One more word for telling what happened.
\end{grammarlens}

\subsection*{Your turn}
\begin{itemize}
\raggedright
  \item \emph{Say it:} \emph{kramśaḥ}
  \item \emph{Say it:} \emph{kramśaḥ}, once more
  \item \emph{Say it:} \emph{nadī kā pānī kramśaḥ ūpar ā rahā thā.} — The river water was rising little by little.
  \item \emph{Recall:} say the Hindi for to swell up, then the Hindi for a dam, then say \emph{kramśaḥ} again
\end{itemize}

\begin{tcolorbox}[breakable,colback=teal!4,colframe=teal!35!black,title={Before you move on}]
What does \hi{क्रमशः} mean? (\textquotedblleft{}Little by little; by degrees\textquotedblright{}.) Say it once more.
\end{tcolorbox}
\section[\texorpdfstring{\hi{सैलाब} (sailāb)}{सैलाब (sailāb)}]{\hi{सैलाब} (sailāb) — a deluge; a flood}

\label{lesson:HI-C305-sailab}



\begin{quote}
Before the new one: say the Hindi for a dam, then the Hindi for little by little.
\end{quote}

\subsection*{You'll want to know: \hi{सैलाब}}
\textbf{\hi{सैलाब}} — \emph{sailāb} — \textquotedblleft{}a deluge; a flood\textquotedblright{}.

\begin{quote}
\hi{इस} \hi{साल} \hi{का} \hi{सैलाब} \hi{पिछले} \hi{साल} \hi{से} \hi{ज़्यादा} \hi{बड़ा} \hi{था।}
\end{quote}

\begin{quote}
\emph{is sāl kā sailāb pichle sāl se zyādā baṛā thā.}
\end{quote}

\begin{quote}
This year's flood was bigger than last year's.
\end{quote}

Masculine.

\begin{cousinweb}
\textbf{\hi{सैलाब}} (\emph{sailāb}) is usually traced to Persian \emph{sailāb}, \textquotedblleft{}a flood\textquotedblright{}, which joins Arabic \emph{sayl}, a torrent, and Persian \emph{āb}, water. Beside \textbf{\hi{बाढ़}}, it is the word for a great rush of water.
\end{cousinweb}

\begin{grammarlens}[title={What you've built}]
One more word for telling what happened.
\end{grammarlens}

\subsection*{Your turn}
\begin{itemize}
\raggedright
  \item \emph{Say it:} \emph{sailāb}
  \item \emph{Say it:} \emph{sailāb}, once more
  \item \emph{Say it:} \emph{is sāl kā sailāb pichle sāl se zyādā baṛā thā.} — This year's flood was bigger than last year's.
  \item \emph{Recall:} say the Hindi for a dam, then the Hindi for little by little, then say \emph{sailāb} again
\end{itemize}

\begin{tcolorbox}[breakable,colback=teal!4,colframe=teal!35!black,title={Before you move on}]
What does \hi{सैलाब} mean? (\textquotedblleft{}A deluge; a flood\textquotedblright{}.) Say it once more.
\end{tcolorbox}
\section[\texorpdfstring{\hi{भँवर} (bhãvar)}{भँवर (bhãvar)}]{\hi{भँवर} (bhãvar) — a whirlpool}

\label{lesson:HI-C305-bhavar}



\begin{quote}
Before the new one: say the Hindi for little by little, then the Hindi for a deluge.
\end{quote}

\subsection*{You'll want to know: \hi{भँवर}}
\textbf{\hi{भँवर}} — \emph{bhãvar} — \textquotedblleft{}a whirlpool\textquotedblright{}.

\begin{quote}
\hi{भँवर} \hi{में} \hi{एक} \hi{नाव} \hi{डूब} \hi{गई।}
\end{quote}

\begin{quote}
\emph{bhãvar meṁ ek nāv ḍūb gaī.}
\end{quote}

\begin{quote}
A boat went down in the whirlpool.
\end{quote}

Masculine.

\begin{cousinweb}
is listed by CDIAL under Sanskrit \emph{bhramara}, \textquotedblleft{}moving unsteadily, revolving\textquotedblright{} (CDIAL 9650), beside \emph{bhramaraka}, a spinning top or a whirlpool, of the family of \emph{bhram}, to whirl about. Water that turns round and round is a \textbf{\hi{भँवर}}.
\end{cousinweb}

\begin{grammarlens}[title={What you've built}]
One more word for telling what happened.
\end{grammarlens}

\subsection*{Your turn}
\begin{itemize}
\raggedright
  \item \emph{Say it:} \emph{bhãvar}
  \item \emph{Say it:} \emph{bhãvar}, once more
  \item \emph{Say it:} \emph{bhãvar meṁ ek nāv ḍūb gaī.} — A boat went down in the whirlpool.
  \item \emph{Recall:} say the Hindi for little by little, then the Hindi for a deluge, then say \emph{bhãvar} again
\end{itemize}

\begin{tcolorbox}[breakable,colback=teal!4,colframe=teal!35!black,title={Before you move on}]
What does \hi{भँवर} mean? (\textquotedblleft{}A whirlpool\textquotedblright{}.) Say it once more.
\end{tcolorbox}
